% chapters/ch10_hare_no_hi.tex
% Chapter 10: Hare-no-Hi OS — Festival Protocol (Complete Integrated Specification)
% Source: NRMO_v5_updated.pdf, Chapter 10 (lines 7091-7269 of v5_layout.txt)

\chapter{Hare-no-Hi OS --- Festival Protocol
(Complete Integrated Specification)}
\label{ch:hare-no-hi}

\nrmobox{Status and scope}{
\textbf{Status}: OPERATIONS --- Context Override Mode\\[0.3em]
This specification assumes NRMO is operating as the daily OS, and
defines the complete operational definition for enjoying Hare-no-Hi
(non-ordinary festival days) without breakdown, safely, and with
full engagement.
}

% =====================================================================
\section{Fundamental Definition (Most Critical)}
\label{sec:hare-definition}

Hare-no-Hi is a non-ordinary festival --- a sacred day.

``Sacred'' does not mean ``correct'' or ``to be revered.'' It means:
a demarcated festival space where different customs from the
ordinary are temporarily permitted.

\nrmobox{Sacredness preservation}{
Sacredness is preserved by \textbf{ending}, \textbf{not deciding},
and \textbf{not carrying over}.
}

The Hare-no-Hi OS does \emph{not} handle ``correctness.'' It is a
state transition that temporarily permits vividness, fluctuation,
excess, and narrative.

% =====================================================================
\section{Relationship with NRMO}
\label{sec:hare-vs-nrmo}

\begin{verbatim}
[ NRMO (Daily OS - Always Running - Ruin Avoidance) ]
                  ^ veto always active
                  |
[ Hare-no-Hi OS (Non-resident - Festival - Time-limited) ]
\end{verbatim}

\begin{itemize}
\item NRMO always retains final veto.
\item Hare-no-Hi OS holds \textbf{no} decision authority whatsoever.
\item Hare does not persuade, overwrite, or evaluate NRMO.
\end{itemize}

% =====================================================================
\section{Purpose}
\label{sec:hare-purpose}

\begin{itemize}
\item Permit narrative.
\item Amplify emotion.
\item Experience the non-ordinary as a festival.
\end{itemize}

\noindent
Judgement, conclusion, and normalisation are \textbf{not} purposes.

% =====================================================================
\section{Permitted Actions (Festival)}
\label{sec:hare-permitted}

\begin{itemize}
\item Exaggerate, dramatise, produce.
\item Play with meaning-making (do not make truth).
\item Optimism, elation, immersion.
\item ``Somehow the best'' --- proceed.
\end{itemize}

% =====================================================================
\section{Prohibitions (Actions That Destroy the Sacred)}
\label{sec:hare-prohibited}

\begin{itemize}
\item Major judgements.
\item Irreversible decisions.
\item Life policy, contracts, promises.
\item Carrying over to the next day and beyond.
\item Invalidating or persuading NRMO.
\end{itemize}

\nrmobox{Sacred but not authority}{
\centering
\textit{Sacredness grants permission to play, not authority to
decide.}
}

% =====================================================================
\section{Activation Conditions}
\label{sec:hare-activation}

\begin{itemize}
\item Can explicitly declare ``Today is Hare.''
\item Start and end times are set.
\item Understood that returning afterward is accepted.
\end{itemize}

\subsection{Activation Checklist}

\begin{itemize}
\item[$\square$] Today's purpose is experience, not judgement.
\item[$\square$] Time frame is clear.
\item[$\square$] Can declare not making major decisions.
\item[$\square$] Understand will not use results tomorrow.
\item[$\square$] Accept returning even when excited.
\item[$\square$] Have accepted NRMO's veto.
\end{itemize}

\noindent
If even one is ambiguous: activation prohibited.

\subsection{Activation Declaration (Ritual)}

\begin{lstlisting}[basicstyle=\ttfamily\small,frame=single]
Today is Hare.
This is a festival.
I will not judge.
I will not decide.
When it ends, I return.
\end{lstlisting}

% =====================================================================
\section{Forced Termination Triggers}
\label{sec:hare-termination}

\begin{itemize}
\item Time exceeded.
\item Fatigue or over-concentration.
\item The moment ``Let's go with this'' appears.
\item Self-identification, sense of mission, conviction.
\end{itemize}

\noindent
Upon detection: immediately return to NRMO. No explanation,
evaluation, or lingering is performed.

% =====================================================================
\section{NRMO Intervention Specification}
\label{sec:hare-intervention}

Upon forced-termination trigger detection:

\begin{itemize}
\item No explanation.
\item No discussion of right or wrong.
\item No handling of emotion.
\item STATE = DAILY\_NRMO.
\end{itemize}

% =====================================================================
\section{Recovery from Misuse Day}
\label{sec:hare-misuse}

\paragraph{Definition of misuse.}
\begin{itemize}
\item Made important judgement during Hare.
\item Tried to carry over results.
\item Broke constraints because ``it was special.''
\end{itemize}

\paragraph{During the same day.}

\begin{lstlisting}[basicstyle=\ttfamily\small,frame=single]
Today had misuse.
No evaluation.
No looking for reasons.
Stop here.
\end{lstlisting}

\begin{itemize}
\item All decisions that day: invalidated.
\item Do not save, share, or reflect.
\item Cut off with body (sleep, disconnect).
\end{itemize}

\paragraph{Next day.}
\begin{itemize}
\item No retrospection.
\item No making lessons.
\item Return to normal NRMO operation.
\item Treat as if it did not happen.
\end{itemize}

% =====================================================================
\section{Paradise Extension (Unlocks)}
\label{sec:hare-paradise}

\begin{itemize}
\item Excessive production / world-building.
\item Creation that does not aim for completion.
\item One-day relationships / immersion.
\end{itemize}

\noindent
Paradise is one day.

% =====================================================================
\section{Morning After: Return to Reality}
\label{sec:hare-return}

\begin{itemize}
\item Normal NRMO SOP restored.
\item No evaluation or lessons from yesterday.
\item ``The festival is over'' as fact only.
\end{itemize}

% =====================================================================
\section{Final Definition}
\label{sec:hare-final}

\nrmobox{Hare-no-Hi OS architecture}{
\centering
\textit{Hare-no-Hi OS is a temporary festival space within the NRMO
system.}\\[0.3em]
\textit{NRMO maintains veto at all times.}\\[0.3em]
\textit{The festival ends and daily life is restored.}\\[0.3em]
\textit{This is the architecture.}
}

% =====================================================================
\section*{Connections to Other Parts}

\begin{itemize}
\item Hare-no-Hi is a Context Override Mode in the Type ZERO mode
hierarchy (Section~\ref{sec:type-zero-modes-c2}). Unlike CORE,
VENTURE, MISSION, and SHUTDOWN, it operates non-resident: it
activates only under explicit declaration and terminates by time
or trigger.
\item The Anti-Contamination requirement
(Section~\ref{sec:ttm-sop-integration}) is the structural reason
Hare-no-Hi outputs cannot carry over: they are training-context
artefacts, not operational decisions.
\item The narrative random generator (Chapter~\ref{ch:hare-no-hi-narrative})
is the content sub-module operating under this OS.
\end{itemize}
